\section{Introduction}

Vision-language-action (VLA) models generate robot actions from images and a language instruction~\cite{brohan2023rt2,kim2024openvla,pi0,pi05}. Yet visual changes that leave the task unchanged can sharply reduce their success: LIBERO-Plus reports that $\pi_0$ falls from 94.2\% on LIBERO to 13.8\% under camera perturbations~\cite{fei2026liberoplus}. Such failures are naturally studied through how much a perturbation changes the policy's features, commands, or motion (its \emph{sensitivity}); the task success it costs is its \emph{fragility}. We show that the two are different things.

The policy's own randomness exposes the gap. VLAs such as $\pi_0$ act in a closed loop: from each observation they generate a sequence of future actions, execute the first five as controls, and replan, dozens of times per episode. Each sequence starts from a random sample, the \emph{policy noise}, so two unperturbed runs from the same state already move differently. Under sensor noise in the images, $\pi_0$'s end effector ends the first five controls 3.5\,mm from where the clean run put it, while a second clean run lands 6.1\,mm away. Yet a second clean run never changes the task outcome here, and the sensor noise costs 14.3 points of success. The converse also occurs: removing a headlight changes $\pi_{0.5}$'s first five commands 5.1 times as much as a dimming matched in image statistics, and both succeed in 77 of 80 runs.

To locate the gap, we follow a perturbation through the policy, image $\to$ features $\to$ commands $\to$ motion $\to$ outcome, and ask at each link how much of its effect carries over. A clean and a perturbed run that start with the same policy noise see the same physical state only at the episode's first policy call; afterwards the two robots move apart. We therefore compare the features and the first five commands sent to the robot at that call, motion after 5 controls (before any replanning) and after 60 (eleven replans later), and success at the end of the episode (Fig.~\ref{fig:model-measurement}). We ask three questions along this path:
\begin{enumerate}
\item \emph{Image to features:} do different perturbations change the features along shared directions? If so, visual robustness can be studied, and intervened on, at one place inside the policy rather than one perturbation at a time.
\item \emph{Features to commands:} does restoring clean values along these directions move the commands toward the clean run's?
\item \emph{Commands to motion and outcome:} do command changes reach the robot, and does the size of the resulting motion tell which runs fail?
\end{enumerate}

We study frozen, LIBERO-finetuned $\pi_{0.5}$~\cite{pi05} and $\pi_0$~\cite{pi0} on LIBERO~\cite{libero} and LIBERO-Plus under four perturbation families, Background, Camera, Light, and Sensor Noise, in about 4,000 simulator rollouts (Table~\ref{tab:study-map}), with a further 6,048 in a supplementary study (Sec.~\ref{sec:restoration}). We call the change a perturbation causes in the features its \emph{response}. Our main findings are:

\textbf{Different perturbations change the features along shared directions.} In $\pi_{0.5}$, 16 directions per layer, fitted to responses, retain 36.5--83.6\% of the response energy on held-out tasks and perturbation families, against at most 10.2\% for 16 directions fitted to clean features or drawn at random (Figs.~\ref{fig:teaser} and~\ref{fig:token-structure}; Sec.~\ref{sec:shared}).

\textbf{Restoring clean values along these directions reduces command error.} As a causal test, we copy clean values back along the shared directions at a middle layer (L9), given a clean observation of the same state. On held-out tasks this \emph{restoration} moves the first five commands toward the clean run's, reducing their error by 24\% under Camera and 5\% under Sensor Noise. The 16 directions achieve most of what replacing the whole layer does. Task success is also higher, by 1.85 and 3.33 points (Figs.~\ref{fig:teaser}b and~\ref{fig:native-restoration}; Sec.~\ref{sec:restoration}).

\textbf{Command changes reach the robot, but their size is not their cost.} The command change shows directly in the robot's motion after 5 controls and is still detectable after 60. Perturbation families that move the robot more tend to cost more success. For individual runs, however, the size of the early motion change is a poor predictor of failure: the AUROC, where 0.5 is chance, is 0.72 in $\pi_{0.5}$ and 0.64 in $\pi_0$ with all families pooled, and 0.62 and 0.55 within a family. Net success also hides lost and gained successes (Fig.~\ref{fig:behavior}; Sec.~\ref{sec:motion}).

These results give research on VLA robustness two things. One is a common place to intervene: perturbations of different kinds act on the same few directions, and restoring them moves the commands. The other is a yardstick: any correction, ours included, should be judged by paired task outcomes against clean reruns.
